From now on, suppose \( R = \KK[x] \) with \( \KK \) field, thus \( R \) is a PID, so previous
results also applies.

\begin{proposition}
	\label{prop:module-linear-map}
	To give an \( R \)-module \( M \) is equivalent to giving a \( \KK \)-vector space \( V \) and a
	\( \KK \)-linear map \( T: V \to V \).
\end{proposition}

\begin{proof}
	If \( M \) is an \( R \)-module, via the inclusion, we have \( \KK \subseteq \KK[x] \), we get an
	operation \( \KK \times M \to M \) that makes \( M \) a \( \KK \)-vector space. We get
	\[
		\begin{aligned}
			T: M & \to M \\
			T(u) & = xu
		\end{aligned}
	\]
	this is \( \KK \)-linear using the fact that \( M \) is a \( \KK[x] \)-module. Conversely: if \( V= \KK \)-vector space
	and \( T: V \to V \) is a \( \KK \)-linear. We define
	\begin{align*}
		R\times V                 & \to V                                                                             \\
		(\sum_{i=0}^n a_i x^i, u) & \mapsto \sum_{i=0}^n a_i T^i(u)                                                   \\
		                          & \text{ where } T^i = \underbrace{T \circ T \circ \cdots \circ T}_{i\text{ times}}
	\end{align*}
	and it is easy to check that
	\begin{itemize}
		\item This makes \( V \) an \( R \)-module.
		\item The 2 maps we defined are mutual inverses of each other.
	\end{itemize}
\end{proof}

\begin{exercise}
	\leavevmode
	\begin{enumerate}
		\item If \( M = (V,T) \), then \( W \subseteq V \) is an \( R \)-module iff
		      \[
			      \begin{cases}
				      W \subseteq V \text{ subspace.} \\
				      T(W) \subseteq W
			      \end{cases}
		      \]
		\item A morphism of \( R \)-modules \( (V,T) \to (V',T') \) is given by \( f: V \to V' \)
		      being \( \KK \)-linear, such that the following diagram is commutative:
		      \[
			      \begin{tikzcd}
				      V \arrow[r, "f"] \arrow[d, "T"'] & V' \arrow[d, "T'"] \\
				      V \arrow[r, "f"]                 & V'
			      \end{tikzcd}
		      \]
		      In particular, we have
		      \[
			      (V,T) \cong (V',T') \iff T \sim T'
		      \]
		      i.e. there exists \( \KK \)-isomorphism \( f: V \to V \), s.t. \( T' = f\circ T \circ f^{-1} \).
	\end{enumerate}
\end{exercise}
Our goal now is to recover some linear algebra results. We will be interested in studying the case
of \( V \) and \( T \), s.t. \( \dim_{\KK}(V) < \infty \). If \( M = (V,T) \), we get: \( M \) is
finitely generated \( R \)-module, since it is finitely generated over \( \KK \), thus apply the
structural theorem, we have
\[
	M \cong \qo{\KK[x]}{(f_1)} \oplus \cdots \oplus \qo{\KK[x]}{(f_r)}
\]
for some  \( f_1,\ldots, f_r \in \KK[x] \) with \( \deg(f_i) \geq 1 \), s.t. \( f_1 \mid f_2 \mid
\cdots \mid f_r \).
\begin{note}
	\leavevmode
	\begin{itemize}
		\item The \textcolor{red}{key point} here is that we cannot have \( \KK[x] \) factors in \( M \) since
		      we have \( \dim_{\KK}(\KK[x]) = \infty \).
		\item If we assume the \( f_i \) are monic, then they are \underline{unique}, monic means
		      the coefficient of top degree term is \( 1 \). In particular
		      \[
			      f = x^n + c_1 x^{n-1} + \cdots + c_n
		      \]
	\end{itemize}
\end{note}

\section{Morphisms between Free Modules via Matrices}
Now give descriptions of morphism betwen \underline{free modules}. Suppose \( R \) is a commutative
ring, \( F, G \) are finitely generated \( R \)-modules and are free modules. Let \( T: F \to G \)
be \( R \)-linear map. We can choose ordered bases
\[
	\begin{aligned}
		\cB_F: & e_1, \ldots, e_n   \\
		\cB_G: & e_1', \ldots, e_m'
	\end{aligned}
\]
Now if \( T (e_j) = \sum_{i=1}^j a_{ij}e_i'\implies \)
\[
	M = M_{\cB_F, \cB_G} (T) = (a_{ij})_{\substack{1 \le i \le m \\ 1 \le j \le n}}
\]
We get bijection as follows:
\[
	\{R\text{-linear maps } F \to G \} \leftrightarrow M_{m,n}(R)
\]
See its behavior under decomposition, say we also have \( H \) to be finitely generated free \( R
\)-module:
\[
	\cB_H: e_1'', \ldots, e_p''
\]
thus
\[
	F \xrightarrow{T} G \xrightarrow{S} H
\]
\( \implies \)
\begin{equation}
	\label{eq:matrix-morphism}
	M_{\cB_F, \cB_H}(T\circ S) = M_{\cB_G, \cB_H}(S) \cdot M_{\cB_F, \cB_G}(T)
\end{equation}
